\documentclass{article}
\usepackage{amsmath}
\usepackage{amssymb}
\usepackage{enumitem}
\usepackage{array}
\usepackage{pifont}

\title{Display Style in Math Mode}
\author{Pete Doe}
\date{\today}

\begin{document}
	\maketitle
	
	\section*{Introduction}
	The canonical example is taken from page 142 of the \textit{TeXBook}, although we've replaced \texttt{\$\$} by \LaTeX's preferred \texttt{\textbackslash[} and \texttt{\textbackslash]}:
	
	\[
	a_0 + \frac{1}{a_1 + \frac{1}{a_2 + \frac{1}{a_3 + \frac{1}{a_4}}}}
	\]
	
	The default typesetting style can be amended by using the \texttt{\textbackslash displaystyle} command:
	
	\[
	a_0 + \frac{1}{\displaystyle a_1 + \frac{1}{\displaystyle a_2 + \frac{1}{\displaystyle a_3 + \frac{1}{\displaystyle a_4}}}}
	\]
	
	Here's another example which demonstrates the effect of \texttt{\textbackslash textstyle}, \texttt{\textbackslash scriptstyle}, and \texttt{\textbackslash scriptscriptstyle}:
	
	\[
	\begin{array}{rcl}
		f(x) &=& \displaystyle \sum_{i=0}^{n} \frac{a_i}{1+x} \\[6pt]
		\textstyle f(x) &=& \textstyle \sum_{i=0}^{n} \frac{a_i}{1+x} \\[6pt]
		\scriptstyle f(x) &=& \scriptstyle \sum_{i=0}^{n} \frac{a_i}{1+x} \\[6pt]
		\scriptscriptstyle f(x) &=& \scriptscriptstyle \sum_{i=0}^{n} \frac{a_i}{1+x}
	\end{array}
	\]
	
	Depending on the value of \(x\), the equation 
	\[
	f(x) = \sum_{i=0}^{n} \frac{a_i}{1+x}
	\]
	may diverge or converge.
	
	\vspace{1cm}
	
	Inline maths elements can be set with a different style: 
	\(f(x) = \displaystyle \frac{1}{1+x}\).
	
	The same is true for display math material:
	
	\[
	\begin{array}{rcl}
		f(x) &=& \displaystyle \sum_{i=0}^{n} \frac{a_i}{1+x} \\[6pt]
		f(x) &=& \textstyle \sum_{i=0}^{n} \frac{a_i}{1+x} \\[6pt]
		f(x) &=& \scriptstyle \sum_{i=0}^{n} \frac{a_i}{1+x} \\[6pt]
		f(x) &=& \scriptscriptstyle \sum_{i=0}^{n} \frac{a_i}{1+x}
	\end{array}
	\]
	
\end{document}